%% =====================================================================
%%  3 Assumptions and Notations
%% =====================================================================
%%  GUIDE
%%   - One row per assumption: statement, justification, quantitative
%%     check. A bare list of assumptions loses the points this section
%%     exists to win.
%%   - Group by theme (geometry / gas / thermal / flow) if you have more
%%     than six; 66 of 74 gold papers have an explicit assumptions section,
%%     only 30 include a notation table.
%%   - Define every symbol once, here or at first use -- not both.
%% =====================================================================
\section{Assumptions and Notations}
\label{sec:assumptions}

The model rests on five assumptions. Each is followed by the estimate that
justifies it, using the nominal case ($V=\SI{628}{m^3}$, $d=\SI{10}{cm}$,
$T_0=\SI{293}{K}$, $p_0=\SI{101.3}{kPa}$), and by the error that would
appear if the assumption were dropped. The error is measured by rerunning
the code with the effect included or removed, not estimated by eye.

\begin{table}[htbp]
  \centering
  \small
  \setlength{\tabcolsep}{4pt}
  \renewcommand{\arraystretch}{1.15}
  \caption{Assumptions and the justification for each. The last column is
  the change in $t_{0.3}$ when the effect is switched on in the code.}
  \label{tab:assumptions}
  \begin{tabular}{@{}l p{0.27\textwidth} p{0.40\textwidth} >{\raggedleft\arraybackslash}p{0.12\textwidth}@{}}
    \toprule
    \# & Assumption & Justification / validity check & Effect \\
    \midrule
    A1 & Ideal gas, $\gamma=1.40$ &
         $p\ll p_c$; $T$ stays in \SI{210}{K}--\SI{293}{K} & $<\SI{2}{\percent}$ \\
    A2 & Quasi-steady orifice flow &
         $L/c_s=\SI{30}{ms}$, i.e.\ $10^{-4}$ of $t_{0.3}$ & $<\SI{1}{\percent}$ \\
    A3 & Sharp-edged hole, $C_d$ fixed &
         $\mathrm{Re}\approx2\times10^{6}$; $C_d$ calibrated in
         Sec.~\ref{sec:disc-validation} & $\pm\SI{6}{\percent}$ \\
    A4 & Well-mixed interior &
         turbulent mixing \SI{10}{s} $\ll t_{0.3}$ & $<\SI{5}{\percent}$ \\
    A5 & Walls at $T_w=\SI{293}{K}$ &
         $\tau_w=\SI{300}{s}$ is comparable to $t_{0.3}$, so the relaxation
         is modelled, not assumed & $-\SI{2}{\percent}$ \\
    \bottomrule
  \end{tabular}
\end{table}

A1 and A2 are the easy ones: the gas stays far from its critical point and
the module empties far more slowly than sound crosses it. A4 is the one
usually asserted without evidence; the estimate above uses the turbulent
mixing time of the cabin air rather than the venting time, which is the
honest comparison. A5 is not really an assumption --- it defines the
thermal layer of Section~\ref{sec:model-thermal}, and its effect is
measured, not bounded.

\begin{table}[htbp]
  \centering
  \small
  \caption{Notation.}
  \label{tab:notation}
  \begin{tabular}{@{}lll@{}}
    \toprule
    Symbol & Meaning & Unit \\
    \midrule
    $p$, $p_0$, $p_a$ & internal, initial, ambient pressure & Pa \\
    $T$, $T_0$, $T_w$  & gas, initial, wall temperature     & K \\
    $V$, $A$, $d$      & module volume, hole area, hole diameter & m$^3$, m$^2$, m \\
    $\gamma$, $R$, $c_s$ & heat-capacity ratio, gas constant, sound speed & --, J\,kg\inv\,K\inv, m\,s\inv \\
    $C_d$              & lumped discharge coefficient       & -- \\
    $u$, $\dot m$      & exit speed, mass flow rate         & m\,s\inv, kg\,s\inv \\
    $\tau$, $\tau_w$   & venting time scale, wall relaxation time & s \\
    $\Pi_T$            & wall-to-venting time ratio, $\tau/\tau_w$ & -- \\
    $t_{0.3}$          & time to reach $0.3\,p_0$           & s \\
    \bottomrule
  \end{tabular}
\end{table}
